\section{File \texttt{input.pmap}}
\subsection{Purpose}
The text file \texttt{input.pmap} specifies how the \swPorosityMapper executable is to be run for generating the porosity maps. 

\subsection{Format description}
\texttt{input.pmap} is a text file. It contains one keyword per line, with format
\begin{lstlisting}
<KEYWORD> <VALUE>
\end{lstlisting}
The order of the keywords does not matter. Empty lines or lines starting with hash '\#' or equal '=' signs are considered as comment lines.\\
Note that
\begin{itemize}
\item The keywords and its value must be separated by at least one space or tabulation characters.
\item The spelling of the keywords is not case sensitive.\\
\item If a given keyword occurs several times in the file, the last read value overwrites the previously specified ones.
\item All physical parameters (e.g. gravitational acceleration, etc.) are in S.I. units.
\end{itemize}

\subsection{List of allowed keywords}

\begin{description}

\item[model] \defaultvalue{none} Set the porosity model used to generate the maps. Allowed values:
\begin{description}
\item[sp] Single Porosity model
\item[dip] Dual-Integral Porosity model
\item[ddp] Depth Dependant Porosity model
\end{description}

\item[input] \defaultvalue{footprint} Set the format expected from the input file \texttt{buildings.txt}. Allowed values:
\begin{description}
\item[footprint] The file specifies the building spatial footprints coordinates
\item[raster] The file consists in a Cartesian grid
\end{description}

\end{description}

\subsection{Example}
\begin{verbatim}
= ligne de commentaire
# symbole pour commentaire : (# =)
=== Porosity Mapper
model	dip
input	footprint
========
\end{verbatim}
